
\begin{result}{B.21}{Proved}{A complete presentation on the sixty geometric core paths}
The fifty-object, sixty-arrow path diagram (B.32)--(B.33)
has a quotient groupoid equivalent to $\Pi_1M_{0,5}$ by the
following twelve relations.  Its comparison functor sends
every arrow to the path (B.32), with its specified relabelling.
The construction below specifies all basepoint transports,
eleven free graph cycles, their five-point braid words, and
an inverse on the five standard generators.

Number cosets of each group in (B.33) by the lexicographically
least permutation in the coset, with the permutations written
as tuples of images of $0,1,2,3,4$.  Number the sixty arrows
$E_1,\ldots,E_{60}$ in that order.  Use the breadth-first
spanning tree from the first bipyramidal vertex, inspecting
incident arrows in their numerical order.  Its non-tree
arrows are
\[
 E_2,E_3,E_6,E_8,E_9,E_{13},E_{19},E_{22},E_{42},E_{45},E_{48}.
                                                               \tag{B.35}
\]
If $t_v$ is the tree path from the root to $v$, the corresponding
free cycles $g_j$ traverse $t_{s(E)}$, then $E$, then
$t_{t(E)}^{-1}$.  Their exact words are the following.  A signed
integer $k$ denotes $E_k$ for $k>0$ and $E_{-k}^{-1}$ for $k<0$;
within each bracket the paths are traversed from left to right.
\begingroup\small
\[
\begin{array}{c|l}
g_1 &[1,-12,41,-44,20,-17,50,-59,35,-26,2,-11,53,-56,33,-28,39,-46,23,-14]\\
g_2 &[1,-12,54,-55,21,-16,3,-10,52,-57,23,-14]\\
g_3 &[1,-12,41,-44,20,-17,6,-7,39,-46,23,-14]\\
g_4 &[1,-12,54,-55,34,-27,8,-5,49,-60,36,-25]\\
g_5 &[1,-12,41,-44,31,-30,9,-4,38,-47,36,-25]\\
g_6 &[1,-12,41,-44,20,-17,50,-59,35,-26,13,-24,58,-51,30,-31,44,-41,12,-1]\\
g_7 &[14,-23,57,-52,29,-32,19,-18,49,-60,36,-25]\\
g_8 &[14,-23,46,-39,28,-33,22,-15,38,-47,36,-25]\\
g_9 &[14,-23,46,-39,28,-33,56,-53,42,-43,32,-29,52,-57,23,-14]\\
g_{10}&[1,-12,41,-44,31,-30,51,-58,45,-40,27,-34,55,-54,12,-1]\\
g_{11}&[1,-12,41,-44,20,-17,50,-59,48,-37,16,-21,55,-54,12,-1].
\end{array}                                                    \tag{B.36}
\]
\endgroup
With the transport to the normalized configuration
$(0,1,10,20,\infty)$ specified in the proof, their images
in the presentation (B.30) are
\[
\begin{array}{c|l@{\qquad}c|l}
g_1&abyz&g_2&z\\
g_3&abz&g_4&z^{-1}y^{-1}z\\
g_5&axz&g_6&z^{-1}x^{-1}z\\
g_7&z^{-1}xz&g_8&b^{-1}z^{-1}y^{-1}\\
g_9&b^{-1}a^{-1}z^{-1}y^{-1}x^{-1}&g_{10}&abxyz\\
g_{11}&a&&.
\end{array}                                                    \tag{B.37}
\]
Define words in these cycles by
\[
 A=g_{11},\qquad B=g_{11}^{-1}g_3g_2^{-1},\qquad
 X=g_2g_6^{-1}g_2^{-1},\qquad
 Y=g_2g_4^{-1}g_2^{-1},\qquad Z=g_2.              \tag{B.38}
\]
The first six defining relations are
\[
\begin{array}{lll}
 g_1=g_3g_4^{-1},&g_5=g_{11}g_2g_6^{-1},&g_7=g_6^{-1},\\
 g_8=g_2g_3^{-1}g_{11}g_4g_2^{-1},&
 g_9=g_2g_3^{-1}g_4g_6g_2^{-1},&
 g_{10}=g_3g_6^{-1}g_4^{-1}.
\end{array}                                                    \tag{B.39}
\]
The remaining six, with $c(U,V)=V^{-1}U^{-1}VU$, are
\[
\begin{array}{lll}
 AXA^{-1}=Z^{-1}XZ,&
 AYA^{-1}=c(X,Z)Yc(X,Z)^{-1},&
 AZA^{-1}=c(X,Z)Z,\\
 BXB^{-1}=X,&
 BYB^{-1}=Z^{-1}YZ,&
 BZB^{-1}=c(Y,Z)Z.
\end{array}                                                    \tag{B.40}
\]
All twelve are based at the root; their conjugates supply
the relation at any other object.  Formulas (B.36),
(B.38)--(B.40) are a complete explicit system, with no
additional unspecified relators.
\end{result}
\begin{proof}
\emph{Actual configuration paths and transports.}
For an arrow coset choose its lexicographically least
representative $p$.  Relabel (B.32) by $p$, so that the
coordinate of label $j$ is $u_j(r)=K_{p^{-1}(j)}(r)$.
Let
\[
 N_{a,b,c}(w)=\frac{(w-a)(b-c)}{(w-c)(b-a)},
                                                               \tag{B.41}
\]
with its usual projective limits when one of the entries
is $\infty$.  The path in the normalized configuration space
is
\[
 Z_p(r)=\bigl(0,1,x_p(r),y_p(r)\bigr),\qquad
 (x_p,y_p)=\bigl(N_{u_0,u_1,u_4}(u_2),
                  N_{u_0,u_1,u_4}(u_3)\bigr).       \tag{B.42}
\]
All sixty rational functions and endpoint configurations
are determined by (B.41)--(B.42).  The endpoint configuration
depends only on the relevant $H_B$ or $H_Q$ coset, because
normalization is invariant under a common projective
transformation.  Thus these are the actual endpoint functors,
not arbitrary compatible maps of coset sets.

For each of the fifty endpoint configurations $(0,1,x,y)$,
choose the following path $T_{x,y}$ from $(0,1,10,20)$.
First move $10$ to $x$ through $N+iN^2$, leaving $0,1,20$
fixed.  Here $N\geq2$ is the least integer for which both
straight segments avoid $\{0,1,20\}$.  Then move $20$ to
$y$ through $N'+i(N')^2$, leaving $0,1,x$ fixed, with the
same least-integer rule.  Omit a motion whose endpoints
coincide.  None of these fifty values of $x$ is $20$, as
substitution in (B.42) verifies.  These choices always
exist: each forbidden collinearity puts a quadratic
condition on the integer $N$, with at most two roots;
a line cannot contain the parabola $N+iN^2$.  The endpoints
themselves avoid the stationary points.  This specifies
all transports without a homotopy-class choice left over.
The based image of an arrow, in traversal order, is
\[
 T_{s(E)}\,Z_p\,T_{t(E)}^{-1}.                     \tag{B.43}
\]
For a core cycle the intermediate transports cancel, leaving
the transport of that actual cycle to the global basepoint.

\emph{Certified straightening.}
Each rational path (B.42) is homotopic, relative endpoints,
to the straight segment between its endpoint configurations.
Here is a uniform proof of the homotopy, also checked by
exact polynomial certificates.  For any pair of its four
points, let $D(r)$ be their nonzero difference and put
$c=3/4$.  The difference formula for the fractional linear
map (B.41) expresses $D(r)$ as a quotient with at most three
original point differences in the numerator and three in
the denominator; limits at infinity only remove factors.
The original differences are constant real or imaginary
multiples, or multiples of $1+ir$ and $1-ir$.
On $1/2\leq r\leq1$ each nonconstant factor changes its
argument from its value at $c$ by at most
\[
 \max\{\arctan(3/4)-\arctan(1/2),
           \arctan(1)-\arctan(3/4)\}
   =\arctan(2/11)<\frac{\pi}{12}.
\]
For the last inequality, $2/11<2-\sqrt3=\tan(\pi/12)$;
squaring the positive quantities proves it.  Hence the
argument of $D(r)/D(c)$ has absolute value less than
$\pi/2$, and
\[
 \operatorname{Re}\bigl(D(r)\overline{D(c)}\bigr)>0.
                                                               \tag{B.44}
\]
The endpoint chord difference has the same strict
half-plane property, since it is a convex combination
of $D(1/\sqrt3)$ and $D(1)$.  Interpolating between the
rational path and its chord consequently never produces
a collision.  This proves the claimed homotopy.

The program \texttt{b\_core\_paths.py} also checks (B.44)
for all six pair differences on each of the sixty paths.
It writes the numerator and denominator of the real
rational function and uses the exact Sturm root count on
$[1/2,1]$ to certify their constant nonzero signs.  All
sixty paths require only one chord.  Thus the subsequent
crossing calculation represents the specified paths,
with a proved homotopy behind the replacement.

\emph{The braid computation.}
Project a complex point to
$\operatorname{Re}(z)+\operatorname{Im}(z)/37$.
All endpoints of the chords and transport segments have
distinct projected positions.  Their coordinates lie in
$\mathbb Q(i,\sqrt3)$.  On each linear segment, every
crossing time is obtained by solving a linear equation
over $\mathbb Q(\sqrt3)$.  Its order and sign are exact:
the sign of $a+b\sqrt3$ is decided from the signs of
$a,b$ and the comparison of $a^2$ with $3b^2$.
At a crossing the two strands are adjacent in the current
projected order.  In the positively oriented perpendicular
coordinate $\operatorname{Im}(z)-\operatorname{Re}(z)/37$,
the positive generator has the previously left strand
below the other strand.  This agrees with the positive
normal half-turn in B.18.  Simultaneous independent
crossings commute; no crossing involving three strands
occurs in the chosen projection.

This procedure gives each of the sixty words (B.43) in
$B_4$.  They are pure, and their linking number on the
stationary strands $0,1$ is zero.  Thus their images in
$\Gamma_{0,5}$ require no further scalar correction.
Combing by B.18 gives words of length at most four in
$a,b,x,y,z$; their braid-word lengths are at most ten.
The program verifies each combed word by the exact Artin
action.  It then computes the spanning tree and the eleven
cycles (B.36), and gives exactly the table (B.37).
The full arrow words, endpoint configurations, tree and
polynomial certificates are in
\texttt{b\_core\_paths.json}; the program is sufficiently
specified above to reproduce each entry.

\emph{Generation and completeness of the relators.}
Substitution of (B.37) in (B.38) gives exactly
$(a,b,x,y,z)$.  This proves surjectivity on the vertex
group and supplies an explicit inverse on its five
standard generators.  The graph has fifty vertices and
sixty edges and is connected, so its free vertex group
has the eleven displayed generators.  Each equation
(B.39) is immediate by substitution in (B.37).
Use these six equations to eliminate
$g_1,g_5,g_7,g_8,g_9,g_{10}$.  The five remaining free
generators are $g_2,g_3,g_4,g_6,g_{11}$.
The change of free basis (B.38) has the explicit inverse
\[
 g_2=Z,\qquad g_3=ABZ,\qquad
 g_4=Z^{-1}Y^{-1}Z,\qquad
 g_6=Z^{-1}X^{-1}Z,\qquad g_{11}=A.
\]
After that change the remaining relations (B.40) are
precisely the complete presentation (B.30), proved in
B.19.  Therefore the comparison on the vertex group is
an isomorphism.  All fifty objects lie in the same component,
and the tree paths give the natural isomorphisms at all
objects.  This proves the asserted equivalence of groupoids.

The independent program \texttt{b\_core\_generation.py}
checks every substitution and every relator in the exact
$F_3\rtimes F_2$ normal form.  As a further direct check
of generation, it folds the eleven base words in $F_2$
while retaining words in the original generators, obtains
the two base generators, and folds the resulting kernel
words in $F_3$.  The latter graph is already the
three-petal rose, so no saturation step is required.
It records these witnesses as well as the shorter inverse
(B.38).  The resulting twelve nonempty relators and their
expanded words in the sixty original edges are retained
in \texttt{b\_core\_generation.json}.  The proof of their
completeness is the explicit elimination just given.
\end{proof}
